\clearpage
\section*{Zusammenfassung} 
Im Rahmen des Moduls "GeoSensorik und Monitoring II" wurden die Messdaten
der Überwachungskampagne 2026 am Schwanderbärgli bei Schwanden bei Brienz
aufbereitet und ausgewertet.Das Gebiet wird aufgrund seiner bekannten
Hangbewegungen seit mehreren Jahrzehnten geodätisch überwacht. Ziel der
vorliegendenArbeit war es, die erhobenen Messdaten für die Netzausgleichung
vorzubereiten, die Genauigkeit und Zuverlässigkeit des Überwachungsnetzes zu
beurteilen und das bestehende Messnetz durch eine gezielte Reduktion von
Beobachtungen zu optimieren. Dabei sollte untersucht werden, inwieweit
sich der Messaufwand für zukünftige Überwachungskampagnen reduzieren lässt,
ohne die Netzqualität wesentlich zu beeinträchtigen. \bigskip

\todo {müsste noch einen Abschnitt für die Literaturrecherche ergänzt werden?}

Zunächst wurden die tachymetrischen Messdaten der Teilgebiete Ägerdi, Brunni,
Brichen und Sperre mit der Software Leica Infinity aufbereitet. Dabei wurden
die Beobachtungen auf Vollständigkeit und Plausibilität überprüft, auffällige
Messsätze ausgeschlossen sowie fehlerhafte oder fehlende Beobachtungen
korrigiert beziehungsweise ergänzt. Zusätzlich wurden die meteorologischen
Korrekturen insbesondere bei längeren Visuren kontrolliert und angepasst.
Der bereinigte Datensatz bildete anschliessend die Grundlage für die
Netzausgleichungen mit GeoSuite (LTOP). \bigskip

Die Auswertung erfolgte schrittweise über einen provisorischen Abriss, eine
robuste und eine weiche Netzlagerung sowie eine gezwängte Ausgleichung.
Dabei wurden die Beobachtungen anhand ihrer Verbesserungen und der
statistischen Kennwerte beurteilt und auffällige Messungen schrittweise
eliminiert. Eine besondere Herausforderung stellte die Höhenlagerung dar,
da beim Höhenfixpunkt 9000 unterschiedliche Höhenangaben vorlagen. Aufgrund
dieser Unstimmigkeiten wurde entschieden, die Höhenlagerung auf die
langstatischen GNSS-Messungen abzustützen. Für die anschliessende
Netzoptimierung wurde die zweite Variante der gezwängten Ausgleichung
verwendet, bei der keine bestimmbare normierte Verbesserung mehr den
festgelegten Grenzwert überschritt. \bigskip

Die Optimierung des Netzes erfolgte in drei aufeinanderfolgenden Schritten.
Dabei wurden hauptsächlich mehrfach bestimmte und stark überbestimmte Visuren
entfernt. Insgesamt konnten 28 \todo{evtl. 30?} Visuren eliminiert werden, wodurch sich die
Anzahl der Lagebeobachtungen von 353 auf 297 und jene der Höhenbeobachtungen
von 177 auf 150 reduzierte. Die mittlere Lagegenauigkeit verschlechterte
sich dabei lediglich von 1.7 mm auf 1.9 mm und die mittlere Höhengenauigkeit
von 2.1 mm auf 2.2 mm. Die Anzahl der unkontrollierten Beobachtungen blieb
unverändert, und die normierten Verbesserungen lagen weiterhin innerhalb des
festgelegten Grenzwerts. \bigskip

Die Resultate zeigen, dass das Überwachungsnetz durch das gezielte Entfernen
redundanter Visuren vereinfacht werden kann, ohne die Genauigkeit und
Zuverlässigkeit gegenüber dem Ausgangsnetz wesentlich zu verschlechtern.
Allerdings bestehen insbesondere im Gebiet Brunni weiterhin Schwächen
hinsichtlich der Zuverlässigkeit. Zudem werden die geforderten
Lagegenauigkeiten bei einzelnen Punkten bereits im Ausgangsnetz nicht
vollständig erreicht. Mit der dritten Optimierungsstufe wurde deshalb eine
sinnvolle Grenze der Beobachtungsreduktion erreicht. Die Ergebnisse bilden
eine Grundlage für die zukünftige Messplanung und zeigen, welche Visuren
ohne wesentliche Beeinträchtigung der Netzqualität weggelassen werden können.

